% ============================================================
% Ejercicio 2
% ============================================================
\section{Ejercicio 2}

\begin{tcolorbox}[enunciado]
\begin{enumerate}[label=(\alph*), leftmargin=*]
    \item Describe un esquema de codificación para ejemplares del problema anterior (en su versión de decisión).
    
    \textbf{Nota:} No es válido usar bibliotecas o herramientas externas que ya lean formatos predefinidos (como JSON, XML, CSV, etc).

    \item Describe e implementa un algoritmo que lea archivos con el formato propuesto en el inciso anterior e imprima en consola los datos del ejemplar.

    El programa debe recibir como entrada (desde consola):
    \begin{itemize}[leftmargin=*]
        \item Nombre del archivo de entrada con el ejemplar a leer.
        
        El ejemplar debe seguir el esquema de codificación propuesto
    \end{itemize}

    Como resultado de la ejecución, el programa deberá producir como salida:
    \begin{itemize}[leftmargin=*]
        \item Imprimir en consola el total de pedidos disponibles
        \item Imprimir en consola la capacidad máxima de la mochila y la ganancia meta
        \item Imprimir dos ejemplos de pedidos disponibles, mostrando su peso y ganancia
    \end{itemize}
\end{enumerate}
\end{tcolorbox}

\subsection{Solución}
% --- Escribe aquí tu solución ---